%% ma: L11-B11-0005
%% lop: 11
%% chuong: 4
%% bai: 11
%% dang: 11.11.2
%% loai: DS
%% muc: [NB, VD, TH, TH]
%% dap_an: "ĐSSS"
%% nguon: "de-cuong-giua-ky-1-toan-11-nam-2026-2027-truong-thpt-hoang-van-thu-ha-noi-2-8.pdf (M08, Đề cương giữa kì 1 Toán 11 THPT Hoàng Văn Thụ, Hà Nội), Đề minh họa, Phần II Câu 1"
%% kien_thuc: "Bài 10 (giao tuyến, giao điểm), Bài 11 (hai đường thẳng song song, chéo nhau, định lí về giao tuyến)"
%% trang_thai: cho-duyet
%% ly_do: "Dạng mới đề xuất 11.11.2"
\begin{ex}
Cho hình chóp $S.ABCD$ có đáy $ABCD$ là hình thang, đáy lớn $CD$. Gọi $I$, $J$ lần lượt là trung điểm của $AD$, $AB$ và $O$ là giao điểm của hai đường chéo hình thang $ABCD$.
\choiceTF{\True Giao tuyến của hai mặt phẳng $(SAC)$ và $(SBD)$ là đường thẳng $SO$}{Gọi $G$ là trọng tâm tam giác $SDB$, mặt phẳng $(IJG)$ cắt đường thẳng $SC$ tại $H$, trong đó $H$ là giao điểm của $JG$ và $SC$}{Giao tuyến của hai mặt phẳng $(SAB)$ và $(SCD)$ không phải đường thẳng đi qua $S$ và song song với $AB$}{Đường thẳng $IJ$ và $SD$ là hai đường thẳng cắt nhau}
\loigiai{a) Đúng: $S$, $O$ là hai điểm chung.

b) Sai: $JG$ và $SC$ nói chung không đồng phẳng (trung điểm của $BD$ không nằm trên $JC$), nên không cắt nhau. Cách tìm $H$: $IJ\parallel BD$ nên $(IJG)\cap(SBD)$ là đường thẳng qua $G$ song song với $BD$, cắt $SO$ tại $L$; gọi $K=IJ\cap AC$ thì $H=KL\cap SC$.

c) Sai: $(SAB)$, $(SCD)$ có điểm chung $S$ và lần lượt chứa $AB\parallel CD$, nên giao tuyến là đường thẳng qua $S$ song song với $AB$.

d) Sai: $IJ$ nằm trong mặt phẳng đáy, $SD$ cắt mặt phẳng đáy chỉ tại $D\notin IJ$; hai đường thẳng không có điểm chung và không song song ($IJ\parallel BD$), nên chéo nhau.}
\end{ex}
